\documentclass[10pt,a4paper]{amsart}
\usepackage{amsmath,amssymb,amsthm,graphicx}
\usepackage{verbatim}
\newenvironment{reference}{\comment}{\endcomment}
%\newenvironment{reference}{}{}
\newenvironment{slogan}{\comment}{\endcomment}
\newenvironment{history}{\comment}{\endcomment}

% For commutative diagrams we use Xy-pic
\usepackage[all]{xy}

% We use 2cell for 2-commutative diagrams.
\xyoption{2cell}
% Theorem environments.
%
\theoremstyle{plain}
\newtheorem{theorem}[subsection]{Theorem}
\newtheorem{proposition}[subsection]{Proposition}
\newtheorem{lemma}[subsection]{Lemma}

\theoremstyle{definition}
\newtheorem{definition}[subsection]{Definition}
\newtheorem{example}[subsection]{Example}
\newtheorem{exercise}[subsection]{Exercise}
\newtheorem{situation}[subsection]{Situation}

\theoremstyle{remark}
\newtheorem{remark}[subsection]{Remark}
\newtheorem{remarks}[subsection]{Remarks}

\numberwithin{equation}{subsection}

% Macros
%
\def\lim{\mathop{\mathrm{lim}}\nolimits}
\def\colim{\mathop{\mathrm{colim}}\nolimits}
\def\Spec{\mathop{\mathrm{Spec}}}
\def\Hom{\mathop{\mathrm{Hom}}\nolimits}
\def\Ext{\mathop{\mathrm{Ext}}\nolimits}
\def\SheafHom{\mathop{\mathcal{H}\!\mathit{om}}\nolimits}
\def\SheafExt{\mathop{\mathcal{E}\!\mathit{xt}}\nolimits}
\def\Sch{\mathit{Sch}}
\def\Mor{\mathop{\mathrm{Mor}}\nolimits}
\def\Ob{\mathop{\mathrm{Ob}}\nolimits}
\def\Sh{\mathop{\mathit{Sh}}\nolimits}
\def\NL{\mathop{N\!L}\nolimits}
\def\CH{\mathop{\mathrm{CH}}\nolimits}
\def\proetale{{pro\text{-}\acute{e}tale}}
\def\etale{{\acute{e}tale}}
\def\QCoh{\mathit{QCoh}}
\def\Ker{\mathop{\mathrm{Ker}}}
\def\Im{\mathop{\mathrm{Im}}}
\def\Coker{\mathop{\mathrm{Coker}}}
\def\Coim{\mathop{\mathrm{Coim}}}

% Boxtimes
%
\DeclareMathSymbol{\boxtimes}{\mathbin}{AMSa}{"02}

%
% Macros for moduli stacks/spaces
%
\def\QCohstack{\mathcal{QC}\!\mathit{oh}}
\def\Cohstack{\mathcal{C}\!\mathit{oh}}
\def\Spacesstack{\mathcal{S}\!\mathit{paces}}
\def\Quotfunctor{\mathrm{Quot}}
\def\Hilbfunctor{\mathrm{Hilb}}
\def\Curvesstack{\mathcal{C}\!\mathit{urves}}
\def\Polarizedstack{\mathcal{P}\!\mathit{olarized}}
\def\Complexesstack{\mathcal{C}\!\mathit{omplexes}}
% \Pic is the operator that assigns to X its picard group, usage \Pic(X)
% \Picardstack_{X/B} denotes the Picard stack of X over B
% \Picardfunctor_{X/B} denotes the Picard functor of X over B
\def\Pic{\mathop{\mathrm{Pic}}\nolimits}
\def\Picardstack{\mathcal{P}\!\mathit{ic}}
\def\Picardfunctor{\mathrm{Pic}}
\def\Deformationcategory{\mathcal{D}\!\mathit{ef}}
\usepackage{cleveref}
\crefname{originalreference}{original reference}{original references}
\Crefname{originalreference}{Original reference}{Original references}
\makeatletter
\newlabel{situation-iso}{{05MH}{source}{}{sourcepointer-1}{}}
\newlabel{situation-iso@cref}{{[originalreference][0][]05MH}{[1][0][]source}{}{}{}}
\makeatother
\crefname{originallemma}{original lemma}{original lemmas}
\Crefname{originallemma}{Original lemma}{Original lemmas}
\makeatletter
\newlabel{spaces-cohomology-lemma-affine-base-change}{{07U8}{source}{}{sourcepointer-2}{}}
\newlabel{spaces-cohomology-lemma-affine-base-change@cref}{{[originallemma][0][]07U8}{[1][0][]source}{}{}{}}
\makeatother
\crefname{originallemma}{original lemma}{original lemmas}
\Crefname{originallemma}{Original lemma}{Original lemmas}
\makeatletter
\newlabel{spaces-morphisms-lemma-i-star-equivalence}{{04CJ}{source}{}{sourcepointer-3}{}}
\newlabel{spaces-morphisms-lemma-i-star-equivalence@cref}{{[originallemma][0][]04CJ}{[1][0][]source}{}{}{}}
\makeatother
\title{Zero and surjectivity functors of a map versus its affine adjoint, with the closed immersion case}
\author{The Stacks Project Authors}
\date{stacks-project:spaces-flat.tex:local-snapshot-20261012; GNU FDL 1.2 or later}
\begin{document}
\maketitle
\noindent\textit{Editorial note.} All delivered mathematical text is the author's own, copied from the pinned Stacks Project chapter file spaces-flat.tex; no statement or proof step was rewritten, reordered or removed.\par\medskip
\noindent\textit{Editorial note.} The counted proof is the author's complete argument: base change stable adjunction between v and u, vanishing equivalence for adjoint pairs, and a unit factorization argument transferring surjectivity through the exact fully faithful direct image of a closed immersion\par\medskip
\noindent\textit{Editorial note.} The label lemma-relate-zero-affine names the counted statement, and the proof environment immediately following it in the source is delivered verbatim as its proof.\par\medskip
\noindent\textit{Editorial note.} The counted statement is carried by the author's own lemma environment in the chapter source.\par\medskip
\noindent\textit{Allowed standard prerequisites (editorial).}\begin{itemize}
\item The situation fixing a map of quasi-coherent modules on an algebraic space over B and the functors on Sch/B attached to a property of the base changed map
\item Commutation of pullbacks with base change, and commutation of the direct image along an affine morphism with base change
\item The adjunction between pullback and direct image for an affine morphism
\item Exactness and full faithfulness of the direct image functor along a closed immersion on quasi-coherent modules
\end{itemize}\medskip
\ifdefined\hypertarget\hypertarget{sourcepointer-1}{}\fi
\noindent\textit{Source reference.} Original reference 05MH: sources/stacks-spaces-flat/pinned-original-source.tex, source lines 1233--1233. Prerequisite pointer invoked at proof line(s) 1437; it names a standard fact of the same chapter that the counted argument uses, and every step of the counted proof built on it is retained verbatim in the delivered body.\par\medskip
\ifdefined\hypertarget\hypertarget{sourcepointer-2}{}\fi
\noindent\textit{Source reference.} Original lemma 07U8: sources/stacks-chow/tags.txt, source lines 9963--9963. Prerequisite pointer invoked at proof line(s) 1455; it names a standard fact of a sibling chapter that the counted argument uses, and every step of the counted proof built on it is retained verbatim in the delivered body.\par\medskip
\ifdefined\hypertarget\hypertarget{sourcepointer-3}{}\fi
\noindent\textit{Source reference.} Original lemma 04CJ: sources/stacks-chow/tags.txt, source lines 5693--5693. Prerequisite pointer invoked at proof line(s) 1467; it names a standard fact of a sibling chapter that the counted argument uses, and every step of the counted proof built on it is retained verbatim in the delivered body.\par\medskip
\begin{lemma}
\label{lemma-relate-zero-affine}
In Situation \ref{situation-iso} suppose given an affine morphism
$i : Z \to X$ and a quasi-coherent $\mathcal{O}_Z$-module $\mathcal{H}$
such that $\mathcal{G} = i_*\mathcal{H}$. Let
$v : i^*\mathcal{F} \to \mathcal{H}$ be the map adjoint to $u$.
Then
\begin{enumerate}
\item $F_{v, zero} = F_{u, zero}$, and
\item if $i$ is a closed immersion, then $F_{v, surj} = F_{u, surj}$.
\end{enumerate}
\end{lemma}

\begin{proof}
Let $T$ be a scheme over $B$. Denote $i_T : Z_T \to X_T$
the base change of $i$ and $\mathcal{H}_T$ the pullback of $\mathcal{H}$
to $Z_T$. Observe that $(i^*\mathcal{F})_T = i_T^*\mathcal{F}_T$
and $i_{T, *}\mathcal{H}_T = (i_*\mathcal{H})_T$.
The first statement follows from commutativity of pullbacks
and the second from Cohomology of Spaces, Lemma
\ref{spaces-cohomology-lemma-affine-base-change}.
Hence we see that $u_T$ and $v_T$ are adjoint maps as well.
Thus $u_T = 0$ if and only if $v_T = 0$. This proves (1).
In case (2) we see that $u_T$ is surjective if and only if
$v_T$ is surjective because $u_T$ factors as
$$
\mathcal{F}_T \to
i_{T, *}i_T^*\mathcal{F}_T \xrightarrow{i_{T, *}v_T} i_{T, *}\mathcal{H}_T
$$
and the fact that $i_{T, *}$ is an exact functor
fully faithfully embedding the category of quasi-coherent modules on
$Z_T$ into the category of quasi-coherent $\mathcal{O}_{X_T}$-modules.
See Morphisms of Spaces, Lemma \ref{spaces-morphisms-lemma-i-star-equivalence}.
\end{proof}

\end{document}
